% !TEX root = main.tex 
\iffalse
\section{Feasibility of Cryptographic Groups} \label{sec-feas}

Here  we want foundational constructions of groups  that meet PGG1/PGG2, CDL, and more,  establishing they are not  outright impossible and no better than using idealized models. 
   %  We will also aim to achieve \emph{knowledge} PGG1/PGG2 by making  knowledge assumptions on the starting group.
%To instantiate a scheme  in an idealized model, we often need to make novel assumptions about the object used to instantiate it, say $\mathsf{SHA256}$ or an elliptic curve group. Unfortunately, it seems far-fetched to give proofs that these objects actually satisfy such assumptions. Instead, a way to lend evidence to that is show hash functions or groups satisfying such assumptions exist based on other, more-standard primitives.  Such a research agenda has been successful for RO model, yielding constructions of hash functions meeting assumptions like perfect one-wayness~\cite{}, non-malleability~\cite{}, UCE~\cite{}, correlation intractability~\cite{}, among others. We aim to develop an analogous research agenda targeting the GGM and other algebraic models.
  An important property of our constructions will be \emph{unique representation}; every  element will be represented by exactly one bitstring. In particular, this is needed for implementing Schnorr signatures.
  
We will use the abstraction of a \emph{group scheme}~\cite{PKC:AgrHof18,EC:AgrHofKas20} but with unique representation. We define a group scheme as $\Gamma = (\Gamma.\Setup,\Gamma.\Val,\Gamma.\Add)$, where:   $\Gamma.\Setup$ on input $\secparam$ outputs group parameters $\pp$,  generator $1_\Gamma$, and order $m\in\N$.   $\Gamma.\Val$ on inputs $\pp,h$ outputs a bit $d$.   $\Gamma.\Add$ on inputs $\pp,h_0,h_1$ outputs $h$ or $\bot$.  We require for all $\lambda \in \N$ and all  $(\pp,1_\Gamma)$ output by $\Gamma.\Setup(1^\lambda)$, 
$\calV_{\pp} : = \{a \in \bits^*~|~\Gamma.\Val(\pp,a)=1\}$ forms a group under $\Gamma.\Add$, generated by $1_\Gamma$, and $m = |\calV_\pp|$.

%Our main question is now:
% Our main question is now:

\subsection{Warm-Up: An ``Unprovable'' Group Scheme}

To instantiate a RO, a popular approach is to use indistinguishability obfuscation (iO)~\cite{DBLP:journals/jacm/BarakGIRSVY12,FOCs:GGHRSW13,JACM:JaiLinSai24} to obfuscate a puncturable PRF~\cite{STOC:SahWat14}. (For intuition, the reader can just think of ``black-box'' obfuscation and a blockcipher like $\mathsf{AES}$.)
We  define an analogous construction for group schemes. % We describe our  construction using a puncturable pseudorandom \emph{injection}~\cite{TCC:BonKimWu17}. However, our construction can be adapted to suitable puncturable PRFs. 
    Let $\calO$ be an iO and  $F \colon \bits^k \times \bits^\lambda \to \bits^{q(\lambda)}$ be a \emph{puncturable pseudorandom injection} (PRI)~\cite{TCC:BonKimWu17} (which, like $\mathsf{AES}$, is invertible given the key) for some polynomial $q$. Our group scheme uses two programs with hardcoded  key $K \in \bits^k$ and $\lambda$-bit prime $p$.
 \vspace{-1mm}
 \begin{center}
 \begin{tabular}{c|c} 
 \begin{minipage}{1in}
\begin{tabbing}
123\=123\=123\=123\=123\=\kill
\textbf{Program} $\mathsf{P}_0[K, p](L_0,L_1)$: \\
\> If $F^{-1}_K(L_b) \in \Z_p$ for $b\in \bits$ then \\
\>\> $L \gets F_K(F^{-1}_K(L_0) + F^{-1}_K(L_1) \bmod p)$ \\
\>\> Return $L$ \\
\> Else return $\bot$
\end{tabbing} \end{minipage}
 &
\begin{minipage}{1in}
\begin{tabbing}
123\=123\=123\=123\=123\=\kill
\textbf{Program} $\mathsf{P}_1[K, p](L)$: \\
\> If $F^{-1}_K(L) \in \Z_p$ return $1$ \\
\> Else return $0$\\\\
\end{tabbing} \end{minipage}
\end{tabular}
\end{center}

\noindent The associated group scheme  $\Gamma_0[\calO,F] = (\Gamma_0.\Setup,\Gamma_0.\Val,\Gamma_0.\Add)$ is defined as follows:
\begin{newitemize}
    \item On input $\secparam$,  $\Gamma_0.\Setup$ chooses $K \getsr \bits^k$ and $\lambda$-bit prime $p$, and sets  $\widehat{\mathsf{P}}_b \getsr \calO(\mathsf{P}_b[K,b])$ for $b \in \bits$. Its output  is  $((\widehat{\mathsf{P}}_0, \widehat{\mathsf{P}}_1),F(1),p)$.
    \item On inputs $(\widehat{\mathsf{P}}_0,  \widehat{\mathsf{P}}_1),h$,  $\Gamma_0.\Val$ outputs $\widehat{\mathsf{P}}_1(h)$.
        \item On inputs $(\widehat{\mathsf{P}}_0,  \widehat{\mathsf{P}}_1),h_0,h_1$,  $\Gamma_0.\Add$ outputs $\widehat{\mathsf{P}}_0(h_0,h_1)$.
\end{newitemize}

As $F_K$ is an \emph{injection},  $\Gamma_0[\calO,F]$ enjoys unique representation.
%As $F$ is length-doubling, it is injective with overwhelming probability, so $\mathsf{P}_0$ is well-defined and  $\Gamma_0[\calO,F]$ enjoys unique representation.
The non-standard notion of puncturable PRI notwithstanding, assuming post-quantum iO  \emph{this group scheme cannot even be shown DL-hard}. The reason is that Shor's algorithm~\cite{FOCS:Shor94} quantumly breaks DL in group schemes with unique representation.  So, a reduction from iO to DL could be used to quantumly break iO. 
% One possibility would be to prove (computational) indifferentiability from a generic group in the  \emph{ideal obfuscation model} (cf.~\cite{C:JLLW23}), but we want standard-model constructions. 


\subsection{Proposed Group Scheme  and its PGG2 Security}
To bypass the above barrier, we will start with a ``base group'' and bootstrap it using iO.  We also dispense with the need for a puncturable PRI and use standard puncturable PRFs. The intuition is that a group element is a DDH-based commitment to an element in $\Z_p$.  %This idea follows~\cite{PKC:AgrHof18,EC:AgrHofKas20}, but we require unique representation. %Also note that recent work has shown barriers when relying \emph{solely} on the underlying group.
%Let $\GG$ be a group of prime-order $p$ generated by $g$. 
 Below,  for $a \in \Z_p$ we use  $[a]_g := g^a$, and we use $\langle a \rangle$ for the encoding of $a$ as a $\lambda$-bit string. Our group scheme uses two programs with hardcoded keys $K_1,K_2,K_3$ and group parameters $\params := (\GG, g, p)$. Let $\calO$ be an iO and let $F \colon \bits^k \times \Z_p \to \Z_p$ and  $F' \colon \bits^k \times \GG^3 \to \bits^{\lambda}$ be puncturable PRFs. (There do exist constructions of punctured PRFs with group domains~\cite{C:AMNYY18}.) Define:
 
 \begin{center}
\begin{minipage}{1in}
\begin{tabbing}
123\=123\=123\=123\=123\=\kill
\textbf{Program} $\mathsf{Q}_0[K_1, K_2, K_3,\params](h_0,h_1)$: \\
\> Parse $h_b$ as $(h_{b,0},h_{b,1},h_{b,2},c_b)$ for $b \in \bits \;;\; x_b \gets c_b \oplus F'_{K_3}(h_{b,0},  h_{b,1},  h_{b,2})$ for $b \in \bits$ \\
\> If $x_b \notin \Z_p$  for some $b\in\bits$ return $\bot$ \\
\> If $(h_{b,0} = [F_{K_1}(x_b) \cdot F_{K_2}(x_b)]_{g})$ and $(h_{b,1} = [F_{K_1}(x_b)]_{g})$ and $(h_{b,2} = [F_{K_2}(x_b)]_{g})$ for $b \in \bits$ \\
\>\> $y \gets x_0 + x_1 \bmod p \; ; \; h_{2,0} \gets [F_{K_1}(y) \cdot F_{K_2}(y)]_{g}$ \;;\; $h_{2,1} \gets [F_{K_1}(y)]_{g}$ \;;\; $h_{2,2} \gets [F_{K_2}(y)]_{g}$ \\
\>\> $c_2 \gets  \langle y \rangle  \oplus F'_{K_3}(h_{2,0}, h_{2,1}, h_{2,2}) \;;\;$ Return $(h_{2,0},h_{2,1},h_{2,2},c_2)$ \\
\> Else return $\bot$\\
\end{tabbing} \end{minipage}

\begin{minipage}{1in}
\begin{tabbing}
123\=123\=123\=123\=123\=\kill
\textbf{Program} $\mathsf{Q}_1[K_1, K_2, K_3,\params](h)$: \\
\> Parse $h$ as $(h_{0},h_{1},h_{2},c) \; ; \; x \gets c \oplus F'_{K_3}(h_{0}, h_{1}, h_{2})$ for $b \in \bits$ \\
\> If  $x \notin \Z_p$ return 0 \\
\> Else if $(h_{0} = [F_{K_1}(x) \cdot F_{K_2}(x)]_{g})$ and $(h_{1} = [F_{K_1}(x)]_{g})$ and $(h_{2} = [F_{K_2}(x)]_{g})$ return 1 \\
\> Else Return 0 \\
\end{tabbing} \end{minipage}
\end{center}
\vspace{-2mm}
\noindent 
Let $\group$ be a group-parameter generator. 
Then we define the associated group scheme  $\Gamma_1[\calO,F,F',\group] = (\Gamma_1.\Setup,\Gamma_1.\Val,\Gamma_1.\Add)$ as (where  for PGG2 we can replace $g$ with $g_0$):
\begin{newitemize}
    \item On input $\secparam$,  $\Gamma_1.\Setup$ chooses $K_1,K_2,K_3 \getsr \bits^k$ and   $(\params := (\GG, g_0, p)) \getsr \group(\secparam)$, $g \getsr \GG$, and sets  $\widehat{\mathsf{Q}}_b \getsr \calO(\mathsf{Q}_b[K_1,K_2,K_2,\params])$ for $b \in \bits$. It outputs $\pp := ((\widehat{\mathsf{Q}}_0, \widehat{\mathsf{Q}}_1),\widehat{g},p)$ where 
    \begin{newmath}
        \widehat{g} \gets (u_0 = [F_{K_2}(1) \cdot F_{K_2}(1)]_{g}, u_1  =  [F_{K_1}(1)]_{g}, u_2 = [F_{K_2}(1)]_{g},  c = F'_{K_2}(u_0,u_1,u_2) \oplus (\langle 1 \rangle \| 0..0)) \;.
    \end{newmath}%
    \item On inputs $\pp,h$,  $\Gamma_0.\Val$ outputs $\widehat{\mathsf{Q}}_1(h)$.
        \item On inputs $\pp,h_0,h_1$,  $\Gamma_0.\Add$ outputs $\widehat{\mathsf{Q}}_0(h_0,h_1)$.
\end{newitemize}
%Note that for PGG2 we can swap $g$ for fixed generator $g_0$.
%\noindent Again, the group scheme has unique representation with overwhleming probability. 
%With overwhelming probability over the coins of $\Gamma_1.\Setup$, the group scheme has unique representation.
 Note that the first three components of an output of $\mathsf{Q}_0$  serves a ``fingerprint'' for $x_0 + x_1 \bmod p$ and hence $\Gamma_1$  enjoys unique representation with overwhelming probability.  

\begin{task}
Show $\Gamma_1$ achieves \emph{$q$-bounded} (in the adversary's $\Exp$-queries)  PGG2. 
\end{task} 

For this, we will build on the strategy of~\cite{AC:BrzMit14a} in the case of $q$-bounded UCEs.  Their idea is to put composable \emph{auxiliary-input point obfuscation} (AIPO)~\cite{C:BitCan10} inside the obfuscated programs in the hybrid games. To leverage AIPO, they puncture the hash key at the points queried by the UCE source. In our context, since the PGG2 source does not get the group parameters, we can puncture them at the points queried by $\advs$. As in~\cite{AC:BrzMit14a},  we will then need to ``pad''  our programs, so the final parameter-length in our construction grows  with $q$. 

Recalling the PGG2 game from Section~\ref{sec-prime}, let us assume the source makes a single $\Exp$-query (and many  $\Op$ queries). Our game chain will start in the PGG2 game. We focus on handling the $\Exp$-query;   let  $x^*$  be the (unpredictable) $\Exp$-query and $(h_0^*,h_1^*,h_2^*,c^*)$ be the response.  First, we will puncture $K_1,K_2$ at $x^*$ and hide $x^*$ in an AIPO put inside of $\widehat{Q}_0$ and $\widehat{Q}_1$. If this AIPO is triggered during execution of $\widehat{Q}_0$ or $\widehat{Q}_1$, they use a hardcoded output of $F$. By  AIPO security, we can change the hardcoded output to random. Then, using DDH in $\group$, we can change $h_0^*$ to random. We can then puncture $K_3$ at $y^* \gets \mathsf{Dlog}(h_0^*)$ and hardcode the output of $F'$ at this input.  Finally, we can switch this hardcoded output to random, because (with overwhelming probability) it does not affect the behavior of \smash{$\widehat{Q}_0,\widehat{Q}_1$}. Finally, the adversary's view is independent of $b$. 
The case $q >1$ will be challenging because more restrictions on the source will come into play.

%\begin{task}
%Eliminate dependence of the padding-length in the above construction on $q$.
%\end{task}

%The PI's prior work~\cite{TCC:BFHO22} shows, roughly, a PGG1 group scheme for masking sources implies a UCE hash function for \emph{split sources}. If we find such an implication also holds for PGG2, eliminating the dependence on $q$ above will yield the first UCE hash functions for split sources with key-length independent of $q$, which is of independent interest.  
%If infeasible, a weaker goal will be to obtain such group schemes that satisfy an analogue of correlated-input hashing~\cite{TCC:GoyOneRao11}, where leakage is restricted to be other hash values.  %Finally, another fascinating problem is:

%\begin{task}
%    Eliminate the use of obfuscation in the above approach. 
%\end{task}
%Before doing so for PGG, we will first do so in the case of UCE. Note that Zhandry~\cite{C:Zhandry16} gives a construction of a 1-query UCE for computationally unpredictable sources via ELFs.

    \subsection{Beyond PGG2 via Extremely Lossy Functions} \label{sec-beyond}

It is unclear how to achieve PGG1 using $\Gamma_1$, because in our PGG2 proof sketch we needed to puncture the group parameters at the source's queries. So:
%    To move beyond PGG2, our idea is as follows:

%Roughly, an ELF is a function whose key is generated in one of two possible modes: injective or a ``tunable'' lossy mode.  We require that for any adversary $\advA$ running in time at most $t$,  we can tune the lossy key to induce a function whose image-size is polynomial in $t$, but can be distinguished by $\advA$ with at most some inverse-polynomial probability.   Constructions are known from exponential-DDH~\cite{C:Zhandry16}.


\begin{task}
    Modify group scheme $\Gamma_1$ into a group scheme $\Gamma_2$ that achieves $q$-bounded  PGG1.
\end{task}
For this, we will compose the punctured PRFs in the construction with ELFs~\cite{C:Zhandry16}, as in  $F \circ \mathsf{ELF}$ for a punctured PRF $F$. Towards showing $q$-query PGG1, let us again consider the case of a single $\Exp$-query $x^*$. Once we switch the ELFs to lossy mode,  we puncture the public parameters at a random point in the ELF's image.  With inverse-polynomial probability, this point is the output of the ELF on $x^*$, so the query response can be switched to random and we can conclude as before. For $q >1$, if our approach doesn't work we will  try to prove \emph{negative} results, \emph{e.g.} that there is no reduction from PGG1 to certain game-based assumptions. This could mean PGG2 ``saves the day'' because it is achievable and still sufficient in some applications.
Then, we will turn to:
%\begin{task}
 %   Extend the above to $q$-query PGG1 for $q >1$.
%\end{task}
\begin{task} \label{task-cdl}
    Build a conversion function $f$ such that  $(\Gamma_2,f)$-CDL holds with a tight reduction. 
\end{task}
This is a special case of Task~\ref{task-inst}, so if we are already successful there we may not need this one. Namely,  here we use a specific group, and the result will imply Schnorr signatures are tightly secure in that group. As in Task~\ref{task-inst}, our candidate $f$ will be an  obfuscated puncturable PRF. %So far, we can show that a CI hash function for sparse relations, achieved by such $f$~\cite{EC:CCRR18}, satisfies a \emph{weak} $(\GG,f)$-CDL for \emph{any} DL-hard $\GG$: Given $h,g \getsr \GG$, output $R \in \Z_p^*$ such that \smash{$h = R \cdot g^{f(R)}$}. 
%So, it is plausible such $f$ can be shown to satisfy $(\GG,f)$-CDL for stronger $\GG$ via  a direct analysis.

Finally, a ``dream result'' would be to eliminate dependence on $q$ and the use of obfuscation. Here we will take inspiration from one-query UCEs without obfuscation in~\cite{C:Zhandry16,EPRINT:CheMao24}. 

\subsection{More Notions and  Algebraic Structures} \label{sec-more}

Finally, we will move to  bilinear group and group action schemes.  Then, we will  study knowledge assumptions, starting with prime-order groups. \emph{E.g.}, does KEA1 imply KEA3~\cite{C:Damgaard91,C:BelPal04}? Does KEA3 imply $q$-PKE~\cite{AC:Groth10a}? What if we add additional assumptions, \emph{e.g.}~iO, and allow the implied assumption to hold in a different group than the starting one? Positive results with unique representation will be challenging here, as previous constructions realizing the AGM~\cite{EC:AgrHofKas20}  use non-unique representation in an essential way. As a first cut, we will try constructing  oracles that give a generic group where KEA1 holds but KEA3  does not. 
The challenge is that it is difficult to say that an adversary breaks a KEA3 --- we must rule out \emph{any} extractor for it. 
%When adding iO, we conjecture we will get positive results (dropping unique representation at first). 

%We may also modify $\Gamma_2$ above if it is difficult to show this.  
%\begin{task}
%    
%\end{task}
%Indeed, there is currently little evidence that there exist groups where OMDL holds in the standard model; in fact, it was only recently shown that OMDL holds in the GGM~\cite{}. To tackle the problem, we will start with OMDL in the novel setting of \emph{group actions}. % At a high level, to achieve such group actions we will obfuscate a program that has keys $K_1,K_2$ hardcoded and on input $x \in \Z_p$ returns $g_0^y$ where $y \gets \mathsf{ELF}_{K_2}(F_{K_1}(x))$ where $\mathsf{ELF}$ is an ELF and $F$ is a puncturable PRF. In the proof, we will puncture a challenge inputs $x_1,\ldots,x_n \in \Z_p$ and, to hide the output points used as exponents and use \emph{multi-bit}-AIPO~\cite{EC:CanDak08} in a similar way as~\cite{AC:MurONeZah22}.  %Because we will need an MB-AIPO for each $x_i$, the  parameters grow with $n$ to account for the corresponding padding of the circuit size.

%\begin{task}
%    Eliminate dependence of the above parameter size on $n$.
%\end{task}

%One approach will be by to try to use the \emph{superfluous padding assumption} (SuPA) from~\cite{} but \emph{without} auxiliary input, which is crucial for plausibility of this assumption. 

%\paragraph*{Eliminating obfuscastion.} Finally, an interesting question is whether we can eliminate obfuscation in any of our constructions.  Indeed, some steps were recently taken by~\cite{} to do this for the UCE construction of~\cite{}  using lattices.   We will explore their techniques in the proposed work.
% This approach was recently taken by the PI in a different context~\cite{}. %Finally:

%\begin{task}
 %   Construct a \emph{group} (rather than group action) where OMDL holds.
%\end{task}
\iffalse
\subsection{The Broader  Picture}

To conclude, we will generalize the above questions to give us better insight into the relative strength of group-based assumptions. For group-based assumptions $\mathsf{A}, \mathsf{B}$, we can ask:
\begin{ques}
 Using  primitive $\mathsf{X}$, can a $\mathsf{B}$-hard group $\GG_{\mathsf{B}}$ be used to build an $\mathsf{A}$-hard group $\GG_{\mathsf{A}}$?
\end{ques}
Above, we aim for a positive result in the case $\mathsf{B}$ is DDH, $\mathsf{A}$ is PGG1/PGG2, and $\mathsf{X}$ is iO.  Clearly, many other choices are possible, \emph{e.g.} $\mathsf{X}$ may be ``nothing.'' Indeed, very basic cases like $\mathsf{B}$ is CDH, $\mathsf{A}$ is DDH, and $\mathsf{X}$ is nothing are surprisingly open to the best of our knowledge. For the weaker case $\GG_{\mathsf{A}} =  \GG_{\mathsf{B}}$, Maurer~\cite{} gives some separations in his version of the GGM.  
\begin{task}
    Develop a  framework for black-box separations addressing the general case above.
\end{task}
\fi
%\paragraph*{Best cryptographic group?}
%Also, in light of our constructions we wonder if there is actually a ``best'' cryptographic group.
%Zhandry studied this question for hash functions~\cite{C:Zhandry22a}, asking whether there is a hash function that for all properties $\mathsf{P}$, satisfies $\mathsf{P}$ if \emph{any} hash function satisfies $\mathsf{P}$.  His answer was ``no,'' but the case of groups is unclear.



%through ELF post-composition.} Recently, the PI demonstrated that \emph{postcomposition} of ELFs with punctured PRFs, as in $F \circ \mathsf{ELF}$, is useful for achieving CCA-secure encryption without random oracles, particularly through common transforms like OAEP and Fujisaki-Okamoto. Here we aim to extend these results to construct group schemes satisfying the \emph{one-more discrete-logarithm assumption} (OMDL), which has a CCA flavor.  In particular:
%begin{task}
%Prove that our basic construction augmented with ELF post-composition satisfies OMDL.
%\end{task}
\iffalse
\subsection{Constructing Post-Quantum ELFs}

To supplement our study, we will aim to construct \emph{post-quantum}  ELFs, which is of independent interest.  In particular, we will look at constructions from iO plus assumptions forms of learning parity with noise (LPN)  that are plausibly both post-quantum and exponentially-secure (the latter being  necessary per~\cite{C:Zhandry16}).  In particular, we will consider LPN with constant noise and linear number of samples.  Recent work constructs (trapdoor but non-extreme) lossy functions from related assumptions~\cite{}, lending plausibility to the approach. 

\fi
\fi 
